\section{Dataset}
\label{sec:dataset}

O projeto utiliza o dataset público \textbf{Synthetic Chess Board Images}
\cite{dataset-chess}, disponível no Kaggle sob licença CC0 (domínio público). Suas
características principais:

\begin{itemize}[nosep]
  \item \textbf{Volume:} aproximadamente 1\,900 imagens de tabuleiros fotografados em
        ângulo (não \textit{top-down}), resolução $1280 \times 1280$ pixels, formato
        JPEG.
  \item \textbf{Material:} tabuleiro e peças de madeira (\textit{boxwood} para peças
        claras, \textit{ebony} para escuras), o que produz baixo contraste entre peça e
        fundo em diversas condições de iluminação.
  \item \textbf{Anotações:} um arquivo JSON por imagem contendo (a) a configuração da
        posição --- \texttt{config} --- mapeando cada casa ocupada ao tipo e cor da
        peça, e (b) os quatro cantos do tabuleiro --- \texttt{corners} --- normalizados
        no intervalo $[0,1]$ (multiplicados pelas dimensões da imagem para obter
        coordenadas em pixels).
  \item \textbf{Superfícies, iluminação e ângulos de câmera variados} entre as imagens,
        o que exige que a detecção de cantos e a classificação de ocupação/cor
        generalizem além de um único cenário fixo.
\end{itemize}

O dataset não contém sequências de jogo contínuas --- cada imagem é uma posição
independente, gerada separadamente. Isso é relevante para a Seção~\ref{sec:fen-jogadas}:
a validação do detector de jogadas é feita com posições controladas (montadas
manualmente) e não com transições reais capturadas entre dois instantes do mesmo
tabuleiro.

\subsection{Download e uso}

O dataset é obtido programaticamente via \texttt{kagglehub} e fica em cache local em
\texttt{\textasciitilde/.cache/kagglehub/}; o projeto não versiona os dados brutos. O
bootstrap de ambiente (local ou Google Colab) e o download automático estão
implementados em \texttt{src/setup.py} do repositório do projeto (ver
Seção~\ref{sec:links}).
